% ===================================================================
%  圖表第 7 號 — 冬天嘅村，春天嘅農莊。
%
%  Traduction, faite en 2026, en cantonais ecrit d'aujourd'hui, en
%  caracteres traditionnels, sur l'ido de texto/io. Regles de la
%  colonne : voir 10-toubiu-01.tex. Deux scenes, deux numerotations :
%  les cles portent c1 et c2.
%
%  LA BASSE-COUR EST L'ENDROIT OU LE CANTONAIS A UNE GRAMMAIRE QUE
%  PEKIN N'A PAS. L'ido marque le sexe et l'age des betes par des
%  suffixes — mutonulo, mutonino, mutonyuno — et le mandarin doit
%  antéposer 公 / 母 / 小. Le cantonais, lui, POSTPOSE, exactement
%  comme l'ido :
%
%    公  le male    雞公 (48), 鴨公 (67), 火雞公 (68), 羊公 (50)
%    乸  la femelle 雞乸 (43), 鴨乸 (66), 羊乸 (53), 豬乸 (77)
%    仔  le petit   雞仔 (75), 羊仔 (54), 豬仔 (78), 兔仔 (79)
%
%  Pekin écrit 公鸡, 母鸡, 小鸡 ; le cantonais écrit 雞公, 雞乸,
%  雞仔. C'est le meme ordre que « hanulo, hanino, hanyuneto », et ce
%  tableau en aligne vingt. 乸 n'existe pas dans l'ecriture du nord.
%
%  LE RESTE DE LA FERME SE SEPARE PAR LE VOCABULAIRE :
%
%    雞竇  le poulailler (Pekin : 鸡舍)   白鴿  le pigeon (Pekin : 鸽子)
%    牛欄  l'etable     (Pekin : 牛棚)    牛油  le beurre (Pekin : 黄油)
%    芝士  le fromage   (Pekin : 奶酪)    河仔  le ruisseau
%    頸巾  le cache-nez (Pekin : 围巾)    鉸剪  les ciseaux (Pekin : 剪刀)
%    飛髮佬 le barbier  (Pekin : 理发师)  花王  le jardinier
%
%  « 竇 » est le mot cantonais du nid et du repaire, et il fait le
%  poulailler comme il fait le trou du rat ; « 芝士 » est l'anglais
%  cheese passe par Hong Kong, la ou Pekin a traduit.
%
%  LES METIERS DU VILLAGE REPRENNENT LE SUFFIXE 佬 DU TABLEAU 6 :
%  打鐵佬 (5), 補鞋佬 (7), 飛髮佬 (8), 造車佬 (9), 造鞍佬 (10),
%  磨粉佬 — le meme procede qui donnait 泥水佬 et 油漆佬 au batiment.
%
%  QUATRE BLOCS ONT DEMANDE L'APPOSITION, pour la raison ordinaire :
%  l'ido nomme la chose avant son possesseur, le chinois l'inverse.
%  Voir c1-10-1 (la glace du ruisseau), c1-17-1 (le tas de neige pres
%  du pont), c2-05-1 (la litiere de la jument) et c2-08-1 (la niche
%  du chien).
% ===================================================================

%%K t07-noto-1 noto
\VUnotes{143.2mm}{%
(*) 想知食飯嘅細節，睇圖表第 6 號同第 13 號。%
}

%%K t07-apar-1 apar
\VUsaut{4.0mm}
\VUcentre{13.2pt}{90}{第二組}
\VUsaut{3.0mm}
\VUfilet{16mm}
\VUsaut{5.6mm}
\VUtitre{15.0pt}{60}{圖表第 7 號}
\VUsaut{3.0mm}

%%K t07-tit-1 sub
\VUcentre{12.0pt}{0}{\textit{第一場。}}
\VUsaut{3.0mm}

%%K t07-tit-2 sub
\VUpk{10.2pt}{40}{冬天嘅村。}
\VUsaut{4.0mm}

%%K t07-c1-01-1 p
1. --- 新年假期嗰陣，我跟住叔叔去咗新城，個細村嚟嘅，佢喺嗰度有幾單生意
要傾。

%%K t07-c1-02-1 p
2. --- 我哋搭嘅係城同嗰條村之間行嘅\VUgras{驛車} \textsuperscript{(11)}；不過
天時凍到痺，坐架咁唔舒服嘅車行一趟，真係唔多好受。

%%K t07-c1-03-1 p
3. --- 所以當我哋見到\VUgras{車伕}喺廣場度勒停架車，就喺\textit{白熊}
\VUgras{客棧} \textsuperscript{(2)} 門口，我哋真係鬆一口氣。\VUgras{掌櫃}
\textsuperscript{(4)} 企喺門檻度等我哋，靜靜哋食緊佢支\VUgras{煙斗。}

%%K t07-c1-03-2 p
佢請我哋入去，我都唔使人請第二次，即刻坐埋去灶膛前面，嗰堆藤枝火燒到
噼啪響。嗰陣我先至笑得出：出面嗰陣削骨嘅\VUgras{北風}吹到塊舊\VUgras{招牌}
\textsuperscript{(3)} 吱吱聲，仲捲起一道道黑色嘅漩渦，將煙囪出嘅\VUgras{煙}
\textsuperscript{(1)} 扯走晒。

%%K t07-c1-04-1 p
4. --- 啱啱夠鐘食晏。我哋冇再等，就將人哋端上嚟嗰餐簡單但係好有味嘅飯
食清光 (*)。

%%K t07-tit-3 sub
\VUpk{10.2pt}{40}{幾家探訪。}
\VUsaut{3.0mm}

%%K t07-c1-05-1 p
5. --- 食完晏，我陪叔叔去見佢啲客——佢做保險嘅。我哋第一個去嘅係
\VUgras{打鐵佬} \textsuperscript{(5)}，佢住喺客棧隔籬。佢正話幫匹馬釘蹄。我幾佩服
佢個夥計嘅臂力：佢用條皮圍裙穩穩噉托住隻馬\VUgras{蹄，}打鐵佬就一錘一錘
將塊\VUgras{蹄鐵} \textsuperscript{(6)} 釘實。

%%K t07-c1-06-1 p
6. --- 跟住我哋去咗村入面嗰個\VUgras{補鞋佬} \textsuperscript{(7)} 度，老雅各布。
佢個招牌係隻靚到不得了嘅\VUgras{長靴，}但係佢肯做嘅，就係幫一對穿咗窿嘅舊鞋
換底。

%%K t07-c1-06-2 p
老人家笑住同我哋講：「喺城入面，我叫做『造鞋師傅』，就冇客。喺呢度，我叫做
『補鞋佬』，做極都做唔完。」

%%K t07-c1-07-1 p
7. --- 佢同我哋講起佢隔籬嗰個\VUgras{飛髮佬} \textsuperscript{(8)}，啲客個個都
唔滿意。佢啲\VUgras{剃刀}成日都崩晒口，佢啲\VUgras{鉸剪}又冇一次剪得好。

村入面淨係得佢一個飛髮佬，所以大家梗係要搵佢剃鬚同飛髮啦。何況佢個人
又孤寒又臭串。

%%K t07-c1-08-1 p
8. --- 好彩其他\VUgras{街坊}唔似佢。譬如話\VUgras{造車佬} \textsuperscript{(9)}
就係個好人，勤力又肯幫手。搵佢整車真係一件樂事，佢做嘅嘢次次都執得妥妥當當。

%%K t07-c1-09-1 p
9. --- 老雅各布對\VUgras{造鞍佬} \textsuperscript{(10)} 都冇乜怨言，趕工嗰陣，
佢仲間中去幫手車\VUgras{馬具}添。

%%K t07-c1-09-2 p
叔叔問起佢頭家：「個個都好。我個孫女喺\VUgras{學校} \textsuperscript{(12)} 讀書。
佢個\VUgras{女先生} \textsuperscript{(13)} 好滿意佢，話佢係個好\VUgras{學生}
\textsuperscript{(14)}。佢就唔似我個衰仔嘞；嗰條友淨係識得玩。個
\VUgras{先生} \textsuperscript{(15)} 教極都教唔入。」

%%K t07-tit-4 sub
\VUsaut{3.0mm}
\VUpk{10.2pt}{40}{村頭村尾。}
\VUsaut{3.0mm}

%%K t07-c1-10-1 p
10. --- 老人家跟住講起村入面嘅新聞。就快起間新學校嘞，因為擺喺鄉公所嗰間
已經細到唔夠用。呢單嘢話都冇咁易，買落磨粉佬個花園嘅一半就得。對嗰個
可憐人嚟講，仲要係一筆好賺。天時凍，\VUgras{磨坊} \textsuperscript{(16)} 嘅
\VUgras{磨盤}再磨唔到麥，因為架\VUgras{水輪} \textsuperscript{(17)} 畀啲\VUgras{冰}
\textsuperscript{(25)} 凍住咗——即係\VUgras{河仔} \textsuperscript{(23)} 結嗰啲冰。

%%K t07-c1-11-1 p
11. --- 「講起起樓，你見過我哋間新\VUgras{教堂} \textsuperscript{(18)} 未呀？
披住層雪，靚到似幅畫噉。啲\VUgras{烏鴉} \textsuperscript{(19)} 認唔返自己個舊鐘樓，
唯有孤伶伶噉企喺\VUgras{墳場} \textsuperscript{(20)} 嗰棵大樹上面。」

%%K t07-c1-12-1 p
12. --- 一講到墳場，補鞋佬就記起啲人——叔叔識嘅，舊年過咗身嗰啲。「好世界呀，
唔知添咗幾多個\VUgras{墳} \textsuperscript{(21)} 喺人哋隔籬，就喺\VUgras{柏樹}
\textsuperscript{(22)} 樹蔭下面！死神連我哋個老公證人都收埋咗。成條村嘅人都去咗
送佢\VUgras{殯。}」

%%K t07-c1-13-1 p
13. --- 「喏，接佢手嗰位喺河仔度溜緊冰。佢啱啱先攞咗對\VUgras{溜冰鞋}
\textsuperscript{(26)} 嚟叫我整，條皮帶斷咗。」

%%K t07-c1-14-1 p
14. --- 「河邊嗰個，係新嚟嘅\VUgras{鄉警} \textsuperscript{(27)}。佢做過憲兵，
村入面啲曳仔見到佢都驚。佢哋一見到佢對\VUgras{綁腿} \textsuperscript{(29)} 同頂
\VUgras{軍帽} \textsuperscript{(28)}，即刻鼠晒。」

%%K t07-c1-15-1 p
15. --- 「哈！嗰個係老約瑟，佢趕住\VUgras{車柴嘅細車} \textsuperscript{(30)}
去麵包鋪。——係喎，」叔叔話，「我認得佢套\VUgras{騾} \textsuperscript{(31)} 車。
我啱好有嘢想同佢講，就趁呢個機會啦。再見嘞，雅各布伯伯。」

%%K t07-c1-16-1 p
16. --- 我哋啱啱行出門口，村入面啲細路就散學嘞，一窩蜂衝去個\VUgras{溜冰坡}
\textsuperscript{(33)}，唔理跌親，唔理\VUgras{跌落嚟} \textsuperscript{(34)} 會唔會
跌斷手腳。其中一個喺造車佬度攞返架\VUgras{雪橇仔} \textsuperscript{(32)}，抱個
同伴坐上去，就一路跑一路拉住走。

%%K t07-c1-17-1 p
17. --- 另外嗰啲就衝上個\VUgras{雪堆} \textsuperscript{(37)}——即係\VUgras{橋}
\textsuperscript{(40)} 旁邊嗰個——開始射昨日砌好嗰個\VUgras{雪人}
\textsuperscript{(39)}，仲同佢戴埋頂帽、插支煙斗、孭埋個書包。

%%K t07-c1-18-1 p
18. --- 佢哋唔點理個冰死人嘅風，又唔點理凍。大半個都冇\VUgras{頸巾}
\textsuperscript{(35)} 圍；打完\VUgras{雪球} \textsuperscript{(38)} 之後要暖返身，
一手\VUgras{熱栗子} \textsuperscript{(42)} 就夠——喺老\VUgras{賣嘢婆}雅各賓娜度買嘅。
佢自己就披住條薄過紙嘅\VUgras{披肩} \textsuperscript{(43)}，凍到抖晒。

%%K t07-tit-5 sub
\VUsaut{3.0mm}
\VUcentre{12.0pt}{0}{\textit{第二場。}}
\VUsaut{3.0mm}

%%K t07-tit-6 sub
\VUpk{10.2pt}{40}{春天嘅農莊。}
\VUsaut{4.0mm}

%%K t07-c2-01-1 p
1. --- 冬天雖然有佢嘅樂趣，我哋仲係打從心底歡迎\VUgras{春天}返嚟。

%%K t07-c2-01-2 p
五月天，黃昏時分嘅農家院子，望落幾咁令人開懷！

%%K t07-c2-02-1 p
2. --- 天邊嘅\VUgras{太陽} \textsuperscript{(1)} 慢慢落緊山，將成片\VUgras{郊野}
\textsuperscript{(3)} 浸喺佢紫紅色嘅\VUgras{光} \textsuperscript{(2)} 裏面。勤力嘅
\VUgras{犁田佬} \textsuperscript{(6)} 趕住犁完自己塊\VUgras{田} \textsuperscript{(4)}。

%%K t07-c2-02-2 p
佢用把\VUgras{犁} \textsuperscript{(5)}——套住匹夠力嘅\VUgras{馬} \textsuperscript{(7)}
——劃出條\VUgras{犁溝} \textsuperscript{(8)}，等\VUgras{播種佬} \textsuperscript{(9)}
一把把噉將\VUgras{種子}撒落去，第日就結出金黃嘅麥穗。

%%K t07-c2-03-1 p
3. --- \VUgras{割草佬} \textsuperscript{(11)} 攞住鐮刀，已經將草地嘅草割晒，翻草嘅
就用\VUgras{乾草} \textsuperscript{(10)} 翻嚟翻去曬乾——用\VUgras{草叉}
\textsuperscript{(12)} 同耙——而家佢哋返緊嚟嘞。

%%K t07-c2-03-2 p
有幾個行喺架牛拉嘅重車前面，孭住自己啲傢俬。另外幾個懶啲嘅，就舒舒服服
攤咗喺啲香噴噴嘅乾草上面。

%%K t07-c2-04-1 p
4. --- 行過嗰陣，佢哋友好噉同\VUgras{花王} \textsuperscript{(16)} 打個招呼。佢喺
\VUgras{果園} \textsuperscript{(13)} 度忙緊，一路剪\VUgras{果樹，}一路接
\VUgras{梨樹} \textsuperscript{(14)} 嘅枝。啲\VUgras{李樹} \textsuperscript{(15)}
開緊花；\VUgras{菜園} \textsuperscript{(17)} 落足咗肥，啲菜、啲椰菜同啲生菜都長得好好。

%%K t07-c2-04-2 p
矮牆上面，一隻靚\VUgras{孔雀} \textsuperscript{(18)} 開緊屏，等人哋讚佢啲毛靚。

%%K t07-tit-7 sub
\VUpk{10.2pt}{40}{農家院子。}
\VUsaut{3.0mm}

%%K t07-c2-05-1 p
5. --- \VUgras{馬房} \textsuperscript{(19)} 啲門開晒，望到入面個草架同堆
\VUgras{草墊} \textsuperscript{(23)}——係嗰匹\VUgras{母馬} \textsuperscript{(21)} 嘅，
\VUgras{馬伕} \textsuperscript{(20)} 啱啱先幫佢卸咗套。隻\VUgras{馬駒}
\textsuperscript{(22)} 對腳長到咁滑稽，一路彈跳一路跟實佢阿媽。

%%K t07-c2-06-1 p
6. --- \VUgras{水井} \textsuperscript{(29)} 附近，啲\VUgras{乳牛} \textsuperscript{(25)}
行去飲水，飲嗰個木\VUgras{水槽} \textsuperscript{(26)}，就係佢哋嘅飲水處。隻
\VUgras{牛仔} \textsuperscript{(24)} 趁機食奶；\VUgras{牛公} \textsuperscript{(27)}
嗅到牧場嗰陣香味，開心到吼起上嚟，威威噉揚起佢對夠曬惡嘅\VUgras{角}
\textsuperscript{(28)}。

%%K t07-c2-07-1 p
7. --- 個個都做緊嘢。\VUgras{工人婆} \textsuperscript{(32)} 用手扯住口井條
\VUgras{繩} \textsuperscript{(31)}，用個\VUgras{水桶} \textsuperscript{(30)} 打水
畀啲牲口飲。\VUgras{穀倉} \textsuperscript{(33)} 隔籬，有個男人喺度耙禾稈；
\VUgras{屋} \textsuperscript{(34)} 門口坐住個老\VUgras{紡紗婆} \textsuperscript{(35)}，
用住架紡車同支\VUgras{紡桿，}好似舊時噉紡緊啲\VUgras{亞麻}或者\VUgras{大麻。}

%%K t07-tit-8 sub
\VUsaut{3.0mm}
\VUpk{10.2pt}{40}{農莊主人。}
\VUsaut{3.0mm}

%%K t07-c2-08-1 p
8. --- \VUgras{農莊主人} \textsuperscript{(36)} 指揮住呢一切，逐樣吩咐啲工人。佢
叫人執返把\VUgras{釘耙} \textsuperscript{(40)} 同把\VUgras{鋤頭} \textsuperscript{(39)}
入去有瓦遮頭嘅地方，唔好擺喺個\VUgras{狗竇} \textsuperscript{(38)} 側邊——即係嗰隻
\VUgras{狗} \textsuperscript{(37)} 瞓嗰個。

%%K t07-c2-09-1 p
9. --- 佢把聲嚇親隻\VUgras{白鴿} \textsuperscript{(41)}，佢拍晒對翼飛返入個
\VUgras{鴿舍} \textsuperscript{(42)}；又嚇親啲\VUgras{雞乸} \textsuperscript{(43)}，
佢哋一窩蜂走返入\VUgras{雞竇} \textsuperscript{(44)}。

%%K t07-c2-10-1 p
10. --- \VUgras{羊欄} \textsuperscript{(48)} 前面，就係老\VUgras{牧羊人} \textsuperscript{(45)} 趕緊佢群\VUgras{羊} \textsuperscript{(49)} 返嚟。
佢手上攞住支從來唔離身嘅\VUgras{牧羊杖} \textsuperscript{(46)}，身上着住件
甩晒色嘅\VUgras{罩衫} \textsuperscript{(47)}，將啲\VUgras{羊公} \textsuperscript{(50)}、
\VUgras{羊乸} \textsuperscript{(53)}、\VUgras{羊仔} \textsuperscript{(54)}、
\VUgras{山羊乸} \textsuperscript{(51)} 同\VUgras{山羊仔} \textsuperscript{(52)} 趕入欄。
佢隻忠心嘅\VUgras{狗} \textsuperscript{(55)} 阿古士行喺最後，一路吠一路催啲落後嘅
行快啲。

%%K t07-c2-11-1 p
11. --- 咁嘈咁郁都攪唔亂隻好\VUgras{奶牛} \textsuperscript{(56)} 嘅心情，佢個
\VUgras{鈴} \textsuperscript{(57)} 叮叮噉響，人哋就喺\VUgras{牛欄} \textsuperscript{(58)}
門口幫佢擠奶。

%%K t07-c2-12-1 p
12. 佢好似明白噉，知道自己要出奶，等\VUgras{農婦} \textsuperscript{(60)} 可以喺個
\VUgras{攪奶桶} \textsuperscript{(61)} 度打\VUgras{牛油，}又或者做啲一流嘅
\VUgras{芝士} \textsuperscript{(64)}——擺喺塊板上面滴住水嗰啲，塊板就架喺個
\VUgras{大木盆} \textsuperscript{(63)} 上面。

%%K t07-c2-13-1 p
13. --- 成間\VUgras{奶房} \textsuperscript{(59)} 畀\VUgras{農場工人}
\textsuperscript{(65)} 執得乾乾淨淨，佢啱啱先喺手上攞住嗰個大桶度洗好啲
\VUgras{奶罐} \textsuperscript{(62)}。

%%K t07-tit-9 sub
\VUsaut{3.0mm}
\VUpk{10.2pt}{40}{家禽場。}
\VUsaut{3.0mm}

%%K t07-c2-14-1 p
14. --- 佢着住對厚厚嘅\VUgras{木屐，}行起上嚟有啲論盡，就嚇到隻\VUgras{火雞公}
\textsuperscript{(68)}、啲\VUgras{鴨乸} \textsuperscript{(66)}、啲\VUgras{鴨公}
\textsuperscript{(67)} 同啲\VUgras{鵝} \textsuperscript{(69)} 周圍走，個個擘大
\VUgras{把嘴} \textsuperscript{(70)} 慌失失噉嗌。

%%K t07-c2-14-2 p
隻\VUgras{雞公} \textsuperscript{(71)} 就好鬥啲，佢豎起個紅色嘅\VUgras{雞冠}
\textsuperscript{(72)}，抖一抖\VUgras{條尾} \textsuperscript{(73)} 嘅毛，然之後
「喔喔啼」噉大大聲叫。

%%K t07-c2-15-1 p
15. --- 一隻\VUgras{雞乸阿媽} \textsuperscript{(74)} 帶住啲\VUgras{糞堆}
\textsuperscript{(81)} 上面搵食嘅\VUgras{雞仔} \textsuperscript{(75)} 周圍啄。隻
\VUgras{豬仔} \textsuperscript{(78)} 就走返去搵佢阿媽，即係我哋喺豬欄隔籬見到
嗰隻大隻\VUgras{豬乸} \textsuperscript{(77)}。

%%K t07-c2-16-1 p
16. --- 喺自己嘅籠入面，啲\VUgras{兔仔} \textsuperscript{(79)} 大啖大啖噉食住
啲嫩\VUgras{草} \textsuperscript{(80)}，係農莊主人個女攞嚟畀佢哋嘅。若安娜好錫佢啲
兔仔；佢日日喺雞竇度執完啲新鮮\VUgras{雞蛋} \textsuperscript{(76)} 之後，就嚟探
佢班小朋友。

\VUsaut{6.0mm}
\VUfilet{22mm}
